%% Cascadilla Proceedings Project paper starter — cascadilla document class.
%%
%% The cascadilla class itself is NOT included in this download: it ships
%% with every modern TeX distribution (and with latextemplates.com's
%% compiler), so bundling a copy would only risk pinning you to an older
%% version than your own system already has.
%%
%% Class documentation: texdoc cascadilla   |   https://ctan.org/pkg/cascadilla
%% Licence: cascadilla.cls is LPPL-1.3, (c) 2008-2012 Max Bane.
%%
%% Used by a number of linguistics conference proceedings (e.g. WCCFL,
%% BUCLD) that publish through the Cascadilla Proceedings Project
%% (https://www.lingref.com/cpp/). The class auto-loads natbib; use
%% \cascadillabibliography{<bibfile>} rather than plain \bibliography, so
%% cascadilla.bst is applied automatically.
\documentclass{cascadilla}

\title{Prosodic Cues to Contrastive Focus in Child-Directed Speech}
\author{Ada Lovelace}
\institution{Institute of Analytical Engines}
\completed{\today}
\copyrightyear{2026}

\begin{document}

\maketitle

\section{Introduction}
Caregivers routinely mark contrastive focus prosodically --- with
pitch accent placement, duration, and pauses --- but it is an open
question how early children can exploit these cues independent of
lexical or syntactic marking. This paper reports an eye-tracking study
testing whether 24-month-olds use prosodic prominence alone to resolve
reference in a two-referent display, following the general approach of
\citet{Shosted2006} and the information-theoretic framing in
\citet{Shannon1948}.

Prior corpus work \citep{Goodman1943} established that contrastive
pitch accents are reliably present in caregiver speech directed at
children under three, but did not test whether children use the cue
online.\footnote{We follow the standard convention of reporting pitch
accent type using ToBI labels throughout.}

\section{Method}
\immedsubsection{Participants}
Twenty-four monolingual English-learning toddlers (mean age 24.3
months, range 22--26) were recruited from a database of families who
had previously expressed interest in participating in developmental
research.

\subsection{Stimuli and procedure}
Each trial presented two novel objects side by side while a recorded
voice produced a sentence naming one of them, either with a neutral
prosodic contour or with a contrastive pitch accent on the target
noun. Table~\ref{tab:conditions} summarises the four experimental
conditions.

\begin{table}
    \centering
    \begin{tabular}{lll}
        \emph{Condition} & \emph{Accent} & \emph{Trials} \\
        Neutral--Left & none & 12 \\
        Neutral--Right & none & 12 \\
        Contrastive--Left & L+H* & 12 \\
        Contrastive--Right & L+H* & 12
    \end{tabular}
    \caption{Experimental conditions, pitch accent type, and trial
    counts per participant.}
    \label{tab:conditions}
\end{table}

\subsection{Analysis}
Looking time to the target object was coded offline from video by two
annotators blind to condition (inter-rater agreement $\kappa = 0.91$).

\section{Results}
Children looked significantly longer at the target object in the
contrastive-accent condition than in the neutral condition, consistent
with early sensitivity to prosodic focus marking independent of
lexical cues.

\section{Discussion and conclusion}
\immedsubsection{Implications for acquisition theory}
These results suggest that prosodic focus marking is available as an
independent referential cue well before children reliably produce
contrastive accents themselves, supporting a comprehension-precedes-
production account of this aspect of information structure.

Future work should test whether the effect generalizes to languages
with different prosodic systems, and whether it interacts with joint
attention cues from the speaker.

\cascadillabibliography{references}

\end{document}
